%!TEX root = ../atomic_jsc.tex
% LTeX: language=en
\section{Equivariant Gr\"{o}bner bases}
%
In this section we give the definition and study basic properties of equivariant Gr\"{o}bner bases of $\monoid$-equivariant ideals in $\poly{\K}{\Indets}$ where $\monoid\actson\Indets$ is a \kl{nicely ordered} monoid action.
To illustrate the definition and the properties,
we use the action $\inc{\OrderAtoms}\actson{\OrderAtoms}$.
\arka{use the same notation in the intro? maybe move it earlier}
Here, $\OrderAtoms$ denotes the set of rational numbers and $\inc{\OrderAtoms}$ denotes the set of strictly monotonic functions from $\OrderAtoms$ to itself.
We use the symbol $\OrderAtoms$ to emphasise the fact that here it is considered as a set with a linear order and not a field.
In the latter case we use the symbol $\Q$.
The action $\inc{\OrderAtoms}\actson{\OrderAtoms}$ is \kl{ordered} by definition,
but clearly not \kl{well ordered}.
It is however \kl{nicely ordered}.
For an intuitive argument,
$x^{k_1}_1\dots x^{k_m}_m\gdivleq[\inc{\OrderAtoms}]y^{\ell_1}_1\dots y^{\ell_n}_n$ for monomials $x^{k_1}_1\dots x^{k_m}_m$ and $y^{\ell_1}_1\dots y^{\ell_n}_n$ in $\poly{\K}{\OrderAtoms}$ with $x_1 < \dots < x_m \in\OrderAtoms$ and $y_1 < \dots y_n \in\OrderAtoms$
(i.e., there exists $\pi\in\inc{\OrderAtoms}$ such that $\pi(x_1)^{k_1}\dots \pi(x_m)^{k_m}$ divides $y^{\ell_1}_1\dots y^{\ell_n}_n$)
if and only if the sequence $k_1\dots k_m$ is smaller than the sequence $\ell_1\dots \ell_n$ in the subsequence ordering (i.e., there exists $1 \leq r_1 < \dots < r_m \leq n$ such that $k_i \leq \ell_{r_i}$ for all $i\in\{1,\dots,m\}$).
Since the subsequence ordering is a \kl{well-quasi-order},
$\gdivleq[\inc{\OrderAtoms}]$ is also so.


In \arka{cite section} we make this intuitive argument formal and show how it can be generalised to prove that some other monoid actions are also \kl{nicely ordered}.
%
\arka{push later}
\begin{lemma}
The relation $\gdivleq[\monoid]$ is a \kl{well-quasi-order} on $\mon{\Indets}$ if and only if there is an equivariant function $\varphi$ from $\mon{\Indets}$ to a well-quasi-order $(P,\leq)$ such that $\monelt[p] \gdivleq \monelt[q]$ if $\varphi(\monelt[p]) \leq \varphi(\monelt[q])$.
\end{lemma}
%
\begin{proof}
($\implies$)
Define a relation $\leq$ in $\setof{\orbit{\monelt[p]}}{\monelt[p]\in\mon{\Indets}}$ as
\[
\orbit{\monelt[p]} \leq \orbit{\monelt[q]}
\quad\text{when}\quad \monelt[p] \gdivleq[\monoid] \monelt[q]
\]
If $\gdivleq[\monoid]$ is a \kl{well-quasi-order} then so is $\leq$.
And the function $\varphi \defined \monelt[p] \mapsto \orbit{\monelt[p]}$ is an equivariant function from $\mon{\Indets}$ to the well-quasi-ordered set $(\setof{\orbit{\monelt[p]}}{\monelt[p]\in\mon{\Indets}},\leq)$ such that $\monelt[p] \gdivleq \monelt[q]$ if and only if $\varphi(\monelt[p]) \leq \varphi(\monelt[q])$.

\noindent($\impliedby$)
Pick an infinite sequence $\monelt[p]_1,\monelt[p]_2,\dots$ of monomials in $\mon{\Indets}$.
There must exist $i < j$ such that $\varphi(\monelt[p]_i) \leq \varphi(\monelt[p]_j)$.
But then $\monelt[p]_i \gdivleq \monelt[p]_j$.
\end{proof}
%
Let $\varleq$ denote the total order on $\Indets$ that is preserved by $\monoid\actson\Indets$.
\arka{$\mon{\Indets}$ hasn't been defined yet}
A monomial $\monelt[p]\in\mon{\Indets}$ is sometimes considered as a function from $\Indets$ to $\N$,
where $\monelt[p](x)$ is the exponent of $x$ in $\monelt[p]$.
For instance, if $\monelt[p] = x^2y$ then $\monelt[p](x) = 2$, $\monelt[p](y) = 1$ and $\monelt[p](z) = 0$ for all $z\neq x,y$.
Using $\varleq$, we define the \intro{co-lexicographic} (colex) ordering on monomials as
follows: $\monelt[n] \intro*\revlexlt \monelt[m]$ if there exists an
indeterminate $x \in \Indets$ such that $\monelt[n](x) < \monelt[m](x)$, and such
that for every $y \in \Indets$, if $x \varlt y$ then $\monelt[n](y) =
\monelt[m](y)$.
For example,
for $x \varlt y \varlt z \in\Indets$ we have
$x^2 y \revlexlt xy^2 \revlexlt z \revlexlt y^2z$.
%
Since $\varleq$ is not assumed to be a well-order and $x \varlt y$ implies $x \revlexlt y$ for $x,y\in\Indets$,
$\revlexleq$ can not be guaranteed to be a well-order.
However, it is still well-behaved with respect to multiplication.
\begin{lemma}
For every $\monelt[p],\monelt[q],\monelt[r]$ we have
\begin{enumerate}
\item $1 \revlexleq \monelt[p]$, and
\item if $\monelt[p] \revlexleq \monelt[q]$ then $\monelt[p]\monelt[r] \revlexleq \monelt[q]\monelt[r]$.
\end{enumerate}
\end{lemma}
%
In particular if $\monelt[n] \divleq \monelt[m]$ then $\monelt[n] \revlexleq \monelt[m]$. 

For a polynomial $p \in \poly{\K}{\X}$,
we define the \intro{leading monomial} $\intro*\lm(p)$ of $p$ as the largest monomial appearing in $p$ with respect to $\revlexleq$,
and the \intro{leading coefficient} $\intro*\lc(p)$ of $p$ as the coefficient of $\lm(p)$ in $p$.
We then define the \intro{leading term} $\intro*\lt(p)$ of $p$ as the product of its \kl{leading monomial} and its \kl{leading coefficient},
and the \intro{characteristic monomial} $\intro*\cm(p)$ of $p$ as the product of its \kl{leading monomial} and all the indeterminates appearing in $p$.
We also define the \intro(monomial){domain} of $\monelt[m]$ as the set $\intro*\dom(\monelt[m])$ of indeterminates $x \in \X$ such that $\monelt[m](x) \neq 0$.
Because the coefficients and monomial in question are highly dependent on the ordering $\varleq$,
we allow ourselves to write $\lm[\Indets](p)$ to highlight the precise ordered set of variables that was used to compute the \kl{leading monomial} of $p$.
\arka{is this needed?}
We extend $\dom$ from monomials to polynomials by defining $\dom(p)$ as the union of the \kl(monomial){domains} of all monomials appearing in $p$.
%
\begin{example}
Say $x \varlt y \varlt z \in \Indets$.
Let $p = x^2y^3 - \frac{1}{2}y^2z\in\poly{\Q}{\Indets}$.
Then $\lm(p) = y^2z$, $\lc(p) = -\frac{1}{2}$ and $\lt(p) = - \frac{1}{2}y^2z$.
\end{example}
%
\begin{remark}
The choice of the co-lexicographic ordering to order monomials is not completely arbitrary.
In \arka{refer to section} we discuss which properties of the co-lexicographic order becomes important for computing equivariant Gr\"{o}bner bases.
In this section we also show that some other well-known ordering of monomials such as the lexicographic ordering does not have these properties.
\end{remark}
%
\section{Equivariant Buchberger's algorithm}
%


\AP A function $f \colon X \to Y$ between two sets $X$ and $Y$ equipped with
actions $\group \actson X$ and $\group \actson Y$ is said to be
\intro(func){equivariant} if for all $\gelem \in \group$ and $x \in X$, we have
$f(\gelem \cdot x) = \gelem \cdot f(x)$. For instance, the
\kl(monomial){domain} of a monomial is an \kl{equivariant function}: if $\gelem
\in \group$, then $\gelem \cdot \dom(\monelt[m]) = \dom(\gelem \cdot
\monelt[m])$. Let us point out that the image of an \kl{orbit finite set} under
an \kl{equivariant function} is \kl{orbit finite}, and that the algorithms that
we will develop in this paper are all \kl(func){equivariant}.
%
\section{Weak Equivariant Gröbner Bases}
\label{sec:weakgb}

\AP In this section we prove that a natural adaptation of \kl{Buchberger's
algorithm} to the equivariant setting computes a \kl{weak equivariant Gröbner
basis} of an \kl{equivariant ideal}. This can be seen as an analysis of the
classical algorithm in the equivariant setting. We assume for the rest of
the section that $\Indets$ is a set of indeterminates equipped with a group
$\group$ acting \kl{effectively oligomorphically} on $\X$, and that $\X$ is
equipped with a total ordering $\varleq$ that is \kl(ord){compatible} with the
action of $\group$. The crucial object of this section is the notion of
\kl{decomposition} of a polynomial with respect to a set $H$.

\begin{definition}
  \label{def:decomposition}
  Let $H$ be a set of polynomials. A \intro{decomposition} of $p$
  with respect to $H$ is given by a finite sequence 
  $\mathfrak{d} \defined \seqof{(a_i, \monelt_i, h_i)}[i \in I]$ such that
  \begin{equation}
    p = \sum_{i \in I} a_i \monelt_i h_i \quad ,
  \end{equation}
  where $a_i \in \K$, $\monelt_i \in \mon{\X}$, and $h_i \in H$ for all $i \in I$.
  The \intro{domain of the decomposition}
  that we write $\intro*\domdec(\mathfrak{d})$ is defined as the union
  of the domains of the polynomials $\monelt_i h_i$ for all $i \in I$.
  The \intro{leading monomial of the decomposition} is defined as
  $
    \intro*\lmdec(\mathfrak{d}) \defined \max(\seqof{\lm(\monelt_i h_i)}[i \in I])
  $.
\end{definition}

Leveraging the notion of decomposition, we can define a weakening of the notion
of \kl{equivariant Gröbner basis}, that essentially mimics the classical notion
of \kl{equivariant Gröbner basis} at the level of \kl{decompositions} instead
of polynomials.

\begin{definition}
  An \kl{equivariant set} $\Basis$ of polynomials is 
  a \intro{weak equivariant Gröbner basis} of an \kl{equivariant ideal}
  $\idl$ if $\IdlGen{\Basis} = \idl$, and if for every polynomial $p \in \idl$,
  and decomposition $\mathfrak{d}$ of $p$ with respect to $\Basis$, there
  exists a decomposition $\mathfrak{d}'$ of $p$ with respect to $\Basis$ such that
  $\domdec(\mathfrak{d}') \subseteq \domdec(\mathfrak{d})$,
  and 
  such that $\lmdec(\mathfrak{d}') = \lm(p)$.
\end{definition}

\AP To compute \kl{weak equivariant Gröbner bases}, we use a rewriting
relation. Given $p,r \in \poly{\K}{\X}$, we write $p \intro*\toeucl{H}
r$ if and only if there exists $q \in H$, $a \in \K$, and $\monelt \in
\mon{\X}$ such that $p = a \monelt q + r$, $\dom(r) \subseteq \dom(p)$, and
$\lm[\Indets](r) \revlexlt \lm[\Indets](p)$. To simplify the
notations, we write $r \intro*\pmonlt p$ to denote $\dom(r) \subseteq
\dom(p)$, and $\lm[\Indets](r) \revlexlt \lm[\Indets](p)$, leaving the
ordered set of indeterminates $\Indets$ implicit.
The relation $\pmonleq$ is extended to decompositions by using 
the analogues of $\dom$ and $\lm$ for decompositions.

\begin{lemma}
  \label{lem:chm}
  The quasi-ordering $\pmonleq$ is \kl(ord){compatible} with the action of $\group$,
  and is well-founded on polynomials, and on \kl{decompositions} of polynomials.
\end{lemma}
\begin{proof}
  The first property is immediate because $\dom$, $\lm$, and $\revlexleq$ are
  compatible with the group action $\group$. 
  The second property follows from the fact that $\revlexlt$ is a total
  well-founded ordering whenever one has fixed finitely many possible 
  indeterminates. In a decreasing sequence, the support of the leading 
  monomials is also decreasing, so that sequence contains only finitely many 
  indeterminates, hence we conclude.
  The same proof works for decompositions.
\end{proof}

\AP As a consequence of \cref{lem:chm}, we know that the rewriting relation
$\toeucl{H}$ is \intro{terminating} for every set $H$. Given a set $H$ of
polynomials, and given a polynomial $p \in \poly{\K}{\X}$, we say that $p$ is
\intro{normalised} with respect to $H$ if there are no transitions $p
\toeucl{H} r$. The set of \intro{remainders} of $p$ with respect to $H$ is
denoted $\intro*\rem{H}{p}$, and is defined as the set of all polynomials $r$
such that $p \toeucl{H}^* r$ and $r$ is \kl{normalised} with respect to $H$.
The following lemma states that $\rem{H}{\cdot}$ is a computable function from
our setting.


\begin{lemma}[name={},restate={lem:normalisation}]
  \label{lem:normalisation}
  Let $H$ be an \kl{orbit finite set} of polynomials, and let $p \in \poly{\K}{\X}$ be a
  polynomial. Then $\rem{H}{p}$ is finite.
  Furthermore, this computation
  is \kl(func){equivariant}. In particular, 
  $\rem{H}{K}$ is a computable \kl{orbit finite} set 
  for every \kl{orbit finite} set $K$ of polynomials.
  \proofref{lem:normalisation}
\end{lemma}


\AP Now that we have a quasi-ordering on polynomials, we prove that given
an \kl{orbit finite} set $H$ of generators, we can compute a \kl{weak
equivariant Gröbner basis}. The computation closely follows the classical
\kl{Buchberger's algorithm}. The main idea is to saturate the set of
generators $H$ to remove some \emph{critical pairs} of the rewriting relation
$\toeucl{H}$. Namely, given two polynomials $p$ and $q$ in $H$, we compute the
set $\intro*\CancelPoly{p}{q}$ of cancellations between $p$ and $q$ as the set of
polynomials of the form $r = \alpha \monelt[n] p + \beta \monelt[m] q$ such
that $\lm(r) < \max(\monelt[n] \lm(p), \monelt[m]\lm(q))$, where $\alpha,\beta
\in \K$, and where $\monelt[n], \monelt[m] \in \mon{\X}$. Let us recall that
given two monomials $\monelt[n], \monelt[m] \in \mon{\X}$, one can compute
$\intro*\lcm(\monelt[n], \monelt[m])$ as the least common multiple of the two
monomials, and that this is an \kl{equivariant operation}.
Using this, we can introduce the \intro{S-polynomial} of two polynomials $p$ and $q$
as in \cref{eq:spoly}.
 \begin{equation}
    \label{eq:spoly}
    \intro*\spoly{p}{q} \defined
    \frac{\lcm(\lm(p), \lm(q))}{\lt(p)} \times p
    - \frac{\lcm(\lm(p), \lm(q))}{\lt(q)} \times q
    \quad .
  \end{equation}



\begin{lemma}[restate=lem:spoly,name={S-polynomials}]
  \label{lem:spoly}
  Let $p$ and $q$ be two polynomials in $\poly{\K}{\X}$.
  All the polynomials in $\CancelPoly{p}{q}$ are obtained by multiplying a monomial
  with
  their \kl{S-polynomial} $\spoly{p}{q}$.
  \proofref{lem:spoly}
\end{lemma}


Remark that the \kl{S-polynomial} is equivariant: if $\gelem \in \group$, then
$\spoly{\gelem \cdot p}{\gelem \cdot q} = \gelem \cdot \spoly{p}{q}$. Given a
set $H$, we write $\intro*\spolyset(H) \defined \bigcup_{p,q \in H}
\rem{H}{\spoly{p}{q}}$. We are now ready to define the saturation algorithm
that computes \kl{weak equivariant Gröbner bases}, described in
\cref{alg:weakgb}.
Let us remark that \cref{alg:weakgb}
is an actual algorithm (\cref{lem:weakgb-computable}) that is
\kl(func){equivariant}.

\begin{algorithm}
  \caption{Computing \kl{weak equivariant Gröbner bases} using the algorithm \intro{weakgb}.}
    \label{alg:weakgb}
    \KwIn{An orbit finite set $H$ of polynomials}
    \KwOut{An orbit finite set $\Basis$ that is a \kl{weak equivariant Gröbner basis} of
      $\EqIdlGen{H}$}
    \Begin{
        $\Basis \gets H$\;
        \Repeat{$\Basis$ stabilizes}{
            $\Basis \gets \Basis \cup \spolyset(\Basis)$\;
        }
        \Return{$\Basis$}\;
    }
\end{algorithm}
%
\begin{lemma}
  \label{lem:weakgb-computable}
  \cref{alg:weakgb} is computable and \kl(func){equivariant},
  and produces an \kl{orbit finite} set $\Basis$ if it terminates.
\end{lemma}
%
\begin{proof}
  Observe that it is enough to show that $\spolyset{\Basis}$ is orbit-finite for every orbit-finite set $\Basis$.
  First, we compute $\Basis^2$, which is an \kl{orbit finite} set of pairs,
  because $\Basis$ is \kl{orbit finite} and $\Indets$ is 
  \kl{effectively oligomorphic}.
  Then, noting that $\spoly{-\ }{-}$ is computable and \kl(func){equivariant},
  we conclude
  that $\bigcup_{p,q \in H}{\spoly{p}{q}}$
  is computable and orbit-finite.
  Now, using \Cref{lem:normalisation}, one can compute the set $\spolyset(\Basis)$, which is also orbit-finite.
  Furthermore, one can decide whether the set $\Basis$ stabilises,
  because the membership of a polynomial $p$ in $\Basis$ is decidable,
  since $\group\actson\Indets$ is \kl{effectively oligomorphic} and $\Basis$ is \kl{orbit finite}.
\end{proof}


Let us now use the semantic assumptions to prove the termination of
\cref{alg:weakgb}
(\cref{lem:weakgb-termination})
and the correctness of the resulting orbit finite set
(\cref{lem:weakgb-correctness}).

\begin{lemma}[name={},restate={lem:weakgb-termination}]
  \label{lem:weakgb-termination}
  Assume that the action $\group \actson \Indets$ is \kl{$\omega$-well-structured}.
  Then, 
  \cref{alg:weakgb} terminates
  on every \kl{orbit finite} set $H$ of polynomials.
  \proofref{lem:weakgb-termination}
\end{lemma}


\begin{lemma}
  \label{lem:weakgb-correctness}
  Assume that $\Basis$ is the output of \cref{alg:weakgb}. Then, it 
  is a \kl{weak equivariant Gröbner basis} of the ideal
  $\EqIdlGen{H}$.
\end{lemma}
\begin{proof}
  It is clear that $\Basis$ is a generating set of $\EqIdlGen{H}$, because
  one only add polynomials that are in the ideal generated by $H$ at every step.

  Let $p \in \EqIdlGen{H}$ be a polynomial,
  and let $\mathfrak{d}$ be a \kl{decomposition} of $p$ with respect to
  $\Basis$, that is, a \kl{decomposition} of the form
  \begin{equation}
    p = \sum_{i \in I} \alpha_i \monelt_i p_i
    \quad .
  \end{equation}
  Where $\alpha_i \in \K$, $p_i \in \Basis$, and $\monelt_i \in \mon{\X}$,
  for all $i \in I$.

  Leveraging \cref{lem:chm}, we know that the ordering $\pmonleq$ is
  well-founded. As a consequence, we can consider a minimal 
  decomposition $\mathfrak{d}'$ of $p$ with respect to $\Basis$ such that $\mathfrak{d}'
  \pmonleq \mathfrak{d}$. We now distinguish two cases, depending on whether
  the leading monomial $\lmdec(\mathfrak{d}')$ of the decomposition $\mathfrak{d}'$ is
  equal to the leading monomial of $p$ or not.

  \begin{description}
    \item[Case 1:] $\lmdec(\mathfrak{d}') = \lm(p)$.
      In this case, we conclude immediately,
      as we also have by assumption $\dom(\mathfrak{d}') \subseteq \dom(\mathfrak{d})$.

    \item[Case 2:] $\lmdec(\mathfrak{d}') \neq \lm(p)$.
      In this case, it must be that the set $J$ of indices such that
      $\monelt[l] \defined \lm(\monelt_i p_i) = \lmdec(\mathfrak{d}')$ is non-empty.
      Let us remark that 
      the sum of leading coefficients 
      of the polynomials in $J$ must vanish: $\sum_{i \in J} \alpha_i \lc(p_i) = 0$.
      As a consequence, the set $J$ has size at least $2$.
      Let us distinguish one element $\star \in J$, and 
      write $J_\star = J \setminus \set{\star}$.
      We conclude that 
      $\alpha_\star = - \sum_{i \in J_\star} \alpha_i \lc(p_i) / \lc(p_\star)$.
      Let us now rewrite $p$ as follows:
      \begin{equation}
        p = \sum_{i \in J_\star} \alpha_i 
        \left(\monelt_i p_i - \frac{\lc(p_i)}{\lc(p_\star)} \monelt_\star p_\star\right)
        + \sum_{i \in I \setminus J} \alpha_i \monelt_i p_i
        \quad .
      \end{equation}
      Now, by definition,
      polynomials $\alpha_i \monelt_i p_i$ for $i \in I \setminus J$ have 
      leading monomials
      strictly smaller than $\monelt[l]$.
      Furthermore,
      the polynomials
      $\monelt_i p_i - \frac{\lc(p_i)}{\lc(p_\star)} \monelt_\star p_\star$ for $i \in J_\star$
      cancel their leading monomials, hence they belong
      to the set $\CancelPoly{p_i}{p_\star}$.
      By \cref{lem:spoly}, we know that these polynomials are obtained by
      multiplying the \kl{S-polynomial} $\spoly{p_i}{p_\star}$ by some monomial.
      Because \cref{alg:weakgb} terminated, we know that 
      $\spoly{p_i}{p_\star} \toeucl{\Basis}^* 0$ by construction.

      By definition of the rewriting relation, we conclude that one can rewrite
      $\spoly{p_i}{p_\star}$ as combination of polynomials in $\Basis$ that
      have smaller or equal leading monomials, and do not introduce new
      indeterminates.

      We conclude that
      the whole sum is composed of polynomials with leading monomials 
      strictly smaller than $\monelt[l]$, and using a subset of the indeterminates
      used in $\mathfrak{d}'$, leading to a contradiction
      because of the minimality of the latter. 
      \qedhere
  \end{description}
\end{proof}

As a consequence of the above lemmas, we can now conclude that the 
\cref{alg:weakgb} computes a \kl{weak equivariant Gröbner basis} of the
ideal $\EqIdlGen{H}$, as stated in \cref{thm:weakgb-comput}.

\begin{theorem}
  \label{thm:weakgb-comput}
  Assume that the action $\group \actson \Indets$
  is \kl{$\omega$-well-structured}
  and satisfies our \kl{computability assumptions}.
  Then, the algorithm $\mathsf{weakgb}$ that takes as input an \kl{orbit finite} set $H$ of generators of an
  \kl{equivariant ideal} $\idl$ and computes a \kl{weak equivariant Gröbner
  basis} $\Basis$ of $\idl$.
\end{theorem}
